\section{Evaluation}
\label{sec:eval}

We ask four questions.
\begin{itemize}
  \item[\rqref{1}] Can \tool read the specifications that real REST APIs
    publish, including their references and composed schemas?
  \item[\rqref{2}] Do the \mode{VALID} values of \tool follow their schema
    more often than values from a simpler generator that knows only types and
    enumerations?
  \item[\rqref{3}] On running APIs, how many operations reach a 2XX response
    with each part of \tool (values, id reuse, order), and how does the
    number of unique 5XX responses change?
  \item[\rqref{4}] What does \tool cost, in time per request and in
    requests spent before the operations are reached?
\end{itemize}

\subsection{Subjects}
\label{sec:subjects}
For \rqref{1} and \rqref{2} we use all \NSpecs OpenAPI specifications that
come with the RESTgym APIs~\cite{restgym}: the ten in the RESTgym
repository and the five APIs that REST League 2027 announced for its
benchmark~\cite{restleague}. Together they have \NOpsTotal operations
(Table~\ref{tab:parser}). For \rqref{3} we need running APIs. We select
APIs with the following criteria:
\begin{enumerate}
  \item it is one of the five known REST League 2027 APIs, so the result
    matters for the competition;
  \item it runs from the released jar in a container that we can restart
    from a fresh state before every run, so runs do not affect each other;
  \item it answers requests without the authentication proxy of RESTgym, so
    our runner can talk to it directly.
\end{enumerate}
Pet-clinic (\NPetclinicOps operations, in-memory database) and
gestao-hospital (\NGestaohospitalOps operations, MongoDB) meet all three.
Flight-search expects a login token that RESTgym adds through a proxy
script, kafka-rest-proxy needs a Kafka broker, and notebook-manager needs a
MySQL server. We did not set these up in the time available.

\subsection{Baselines}
\label{sec:baselines}
We compare four configurations. Each adds one part of \tool to the previous
one, so the difference between two neighbours is the effect of that part.
\begin{description}
  \item[Baseline.] Random values that follow only the schema \texttt{type}
    and \texttt{enum}, such as eight random letters for a string, an integer
    from 0 to 1{,}000, and one item per array. This is what a simple
    schema-based random tester sends. It sees the same parsed operations as
    \tool and fills the same parameters.
  \item[+Values.] The value generator of \tool (Step~3), with no id reuse.
  \item[+Pool.] +Values, and ids reused through the resource pool (Step~4).
  \item[+Order (full).] +Pool, and the operation order of Step~2. All other
    configurations use the order of the specification file.
\end{description}
All four use the same runner, the same budget, the same seeds, and a freshly
started API. The runner and the pool are small stand-ins for the team's
engine (Section~\ref{sec:threats}), the same for every configuration.

\subsection{Measurements}
\begin{description}
  \item[\rqref{1}.] Per specification: operations found, compared with the
    method entries in the file; \texttt{\$ref} entries left after parsing;
    median parse time of five runs.
  \item[\rqref{2}.] We generate 50 \mode{VALID} requests per operation with
    each generator and fixed seeds. An \emph{input} is one sent path or query
    parameter, or one request body. An input \emph{conforms} when it
    satisfies every keyword of its schema that our checker understands:
    \texttt{type}, \texttt{enum}, \texttt{required}, \texttt{pattern},
    length and range limits, and the formats \texttt{date},
    \texttt{date-time}, \texttt{email} and \texttt{uuid}.
  \item[\rqref{3}.] Per run: $|\Reached|$ (Equation~\ref{eq:reached}) and
    $|\Faults|$ (Equation~\ref{eq:faults}). We report the mean and standard
    deviation over \NReps runs, and the union of \Faults over all runs.
  \item[\rqref{4}.] Mean time to generate one request; and the requests the
    runner spent in its nominal phase, which ends when every operation was
    either reached or tried ten times.
\end{description}

\subsection{Setup}
\label{sec:setup}
All runs used one laptop (Intel Core 7 240H, 16 logical CPUs, 8~GB of memory
for Docker, Linux 6.18 under WSL2), Python 3.14.4 and Docker 29.8.2. The APIs
ran in the \texttt{eclipse-temurin:17-jre} image from the RESTgym API
repositories at commits \texttt{e8250db} (pet-clinic) and \texttt{d4cb6c7}
(gestao-hospital); gestao-hospital used \texttt{mongo:6} with the seed data
of its repository. Every run had a budget of 20 requests per operation
(\NPetclinicBudget for pet-clinic, \NGestaohospitalBudget for
gestao-hospital). The runner first makes two nominal passes, giving each
operation up to five \mode{VALID} requests per pass until its first 2XX. The
rest of the budget goes to error requests, round-robin over the operations:
\mode{BOUNDARY} and \mode{INVALID} in turn for \tool, random values for the
baseline. Each configuration ran \NReps times, with seeds 1000 to 1004.

\subsection{\rqref{1}: Reading Real Specifications}
\label{sec:rq1}

%% Generated by src/eval/summarize.py from data/parser.csv and data/conformance.csv. Do not edit.
\begin{table}[t]
  \caption{The parser reads all 15 specifications and leaves no reference unresolved; the value generator keeps every VALID input inside its schema. \emph{Ops}: operations. \emph{\$refs}, \emph{allOf}: occurrences in the specification file. \emph{Left}: \$refs left after parsing. \emph{ms}: median parse time of five runs. \emph{Baseline}, \emph{\tool}: share of generated VALID parameter values and bodies that satisfy their schema (50 requests per operation).}
  \label{tab:parser}
  \footnotesize
  \begin{tabular}{lrrrrrrr}
    \toprule
    API & Ops & \$refs & allOf & Left & ms & Baseline & \tool \\
    \midrule
    features-service & 18 & 7 & 0 & 0 & 0.3 & 100.0\% & 100.0\% \\
    flight-search & 40 & 87 & 0 & 0 & 2.1 & 86.4\% & 100.0\% \\
    genome-nexus & 23 & 95 & 0 & 0 & 4.4 & 100.0\% & 100.0\% \\
    gestao-hospital & 20 & 25 & 0 & 0 & 0.7 & 94.3\% & 100.0\% \\
    kafka-rest-proxy & 50 & 588 & 32 & 0 & 5.1 & 98.7\% & 100.0\% \\
    languagetool & 2 & 0 & 0 & 0 & 0.1 & 100.0\% & 100.0\% \\
    market & 13 & 28 & 0 & 0 & 0.6 & 73.9\% & 100.0\% \\
    ncs & 6 & 6 & 0 & 0 & 0.7 & 100.0\% & 100.0\% \\
    notebook-manager & 5 & 6 & 0 & 0 & 0.2 & 71.4\% & 100.0\% \\
    person-controller & 12 & 16 & 0 & 0 & 0.4 & 50.0\% & 100.0\% \\
    pet-clinic & 35 & 149 & 5 & 0 & 2.8 & 80.3\% & 100.0\% \\
    project-tracking-system & 59 & 83 & 0 & 0 & 1.9 & 100.0\% & 100.0\% \\
    restcountries & 22 & 0 & 0 & 0 & 0.3 & 100.0\% & 100.0\% \\
    scs & 11 & 0 & 0 & 0 & 0.3 & 100.0\% & 100.0\% \\
    user-management & 22 & 29 & 0 & 0 & 0.8 & 88.9\% & 100.0\% \\
    \midrule
    total & 338 & 1{,}119 & 37 & 0 & & 94.2\% & 100.0\% \\
    \bottomrule
  \end{tabular}
\end{table}


Table~\ref{tab:parser} shows that \tool found every operation in every
specification and left no reference unresolved. The specifications use
references heavily: \NRefsTotal \texttt{\$ref} entries in total, and
\NRefsMax in \NRefsMaxApi alone, which also has most of the \NAllofTotal
\texttt{allOf} schemas.

\begin{observation}
\label{obs:parse}
\tool parsed all \NSpecs specifications, with \NOpsTotal operations and
\NRefsTotal references, and left \NRefsLeft references unresolved.
\end{observation}

The set has a gap. None of the \NSpecs specifications contains a recursive
schema (we counted \NRecursiveTotal), so on real data the cycle rule of
Step~1 never fired. It is covered only by the small test specifications.

\subsection{\rqref{2}: Values That Follow the Schema}
\label{sec:rq2}
The last two columns of Table~\ref{tab:parser} compare the two generators.
\tool kept \NConfTool of its \NConfToolInputs inputs inside their schema.
The baseline kept \NConfBaseline, and fell below 100\% in
\NConfBaselineBelow of the \NSpecs specifications, down to
\NConfBaselineMin for \NConfBaselineWorstApi. The baseline's \NViolBaseline
violations were almost all of two kinds that type-only generation cannot
meet: \NViolBaselineFormatPattern of them were a \texttt{format} (mostly
\texttt{date} and \texttt{date-time}) or a \texttt{pattern}.

\begin{observation}
\label{obs:conform}
\tool's \mode{VALID} inputs followed their schema in \NConfTool of cases,
against \NConfBaseline for the type-only baseline;
\NViolBaselineFormatPattern of the baseline's violations were formats or
patterns.
\end{observation}

\subsection{\rqref{3}: Operations Reached on Running APIs}
\label{sec:rq3}

%% Generated by src/eval/summarize.py from data/live/. Do not edit.
\begin{table}[t]
  \caption{Each part of \tool adds operations that reach a 2XX response. Mean over 5 runs, each on a freshly started API with the same budget of 20 requests per operation. \emph{Ops 2XX}: operations with at least one 2XX response ($\pm$ standard deviation). \emph{Share}: Ops 2XX divided by all operations. \emph{Nominal}: requests spent in the nominal phase (fewer is better). \emph{5XX}: unique 5XX responses per run. \emph{Union}: unique 5XX over all 5 runs.}
  \label{tab:live}
  \footnotesize
  \begin{tabular}{llrrrrr}
    \toprule
    API & Configuration & Ops 2XX & Share & Nominal & 5XX & Union \\
    \midrule
    \multirow{4}{*}{gestao-hospital} & Baseline & 1.0 $\pm$ 0.0 & 5.0\% & 191.0 & 1.0 $\pm$ 0.0 & 1 \\
     & +Values & 14.2 $\pm$ 1.2 & 71.0\% & 99.4 & 1.6 $\pm$ 0.5 & 2 \\
     & +Pool & 15.4 $\pm$ 0.5 & 77.0\% & 76.0 & 1.4 $\pm$ 0.5 & 2 \\
     & +Order (full) & 15.6 $\pm$ 0.5 & 78.0\% & 75.0 & 1.6 $\pm$ 0.5 & 2 \\
    \midrule
    \multirow{4}{*}{pet-clinic} & Baseline & 10.6 $\pm$ 0.8 & 30.3\% & 276.6 & 9.0 $\pm$ 0.0 & 9 \\
     & +Values & 20.0 $\pm$ 0.9 & 57.1\% & 173.8 & 83.4 $\pm$ 2.9 & 184 \\
     & +Pool & 34.0 $\pm$ 0.0 & 97.1\% & 53.2 & 98.4 $\pm$ 3.3 & 205 \\
     & +Order (full) & 33.4 $\pm$ 0.5 & 95.4\% & 52.6 & 97.8 $\pm$ 3.6 & 194 \\
    \bottomrule
  \end{tabular}
\end{table}


Table~\ref{tab:live} shows the four configurations. On pet-clinic, the
baseline reached \NPetclinicBaselineOps of \NPetclinicOps operations on
average (\NPetclinicBaselineCov). The value generator alone raised this to
\NPetclinicValuesOps (\NPetclinicValuesCov), and id reuse to
\NPetclinicValuesPlusPoolOps (\NPetclinicValuesPlusPoolCov). On
gestao-hospital, the baseline reached \NGestaohospitalBaselineOps of
\NGestaohospitalOps operations (\NGestaohospitalBaselineCov), +Values
\NGestaohospitalValuesOps (\NGestaohospitalValuesCov), and +Pool
\NGestaohospitalValuesPlusPoolOps (\NGestaohospitalValuesPlusPoolCov).

\begin{observation}
\label{obs:values}
Schema-aware values alone raised the operations reached from
\NPetclinicBaselineCov to \NPetclinicValuesCov on pet-clinic and from
\NGestaohospitalBaselineCov to \NGestaohospitalValuesCov on
gestao-hospital.
\end{observation}

\begin{observation}
\label{obs:pool}
Reusing ids from earlier responses raised the operations reached further, to
\NPetclinicValuesPlusPoolCov on pet-clinic and
\NGestaohospitalValuesPlusPoolCov on gestao-hospital.
\end{observation}

The operation order of Step~2 changed little. The full configuration reached
\NPetclinicFullOps operations on pet-clinic against
\NPetclinicValuesPlusPoolOps for +Pool, and \NGestaohospitalFullOps on
gestao-hospital against \NGestaohospitalValuesPlusPoolOps. It spent about
the same number of nominal requests (\NPetclinicFullNominal against
\NPetclinicValuesPlusPoolNominal on pet-clinic, \NGestaohospitalFullNominal
against \NGestaohospitalValuesPlusPoolNominal on gestao-hospital).
Section~\ref{sec:discussion} explains why.

\begin{observation}
\label{obs:order}
Once ids are reused, dependency order changed the operations reached by less
than one operation on average: \NPetclinicValuesPlusPoolOps to
\NPetclinicFullOps on pet-clinic and \NGestaohospitalValuesPlusPoolOps to
\NGestaohospitalFullOps on gestao-hospital.
\end{observation}

\paragraph{The motivating operation.} Over the \NReps baseline runs on
pet-clinic, the runner sent \NMotivBaselineRequests requests to
\texttt{POST /owners/\{ownerId\}/pets}. All \NMotivBaselineServer returned
500 and none created a pet. The full configuration created a pet in
\NMotivFullRunsOk of \NReps runs.

\paragraph{Unique 5XX responses.} The configurations with \tool found far
more unique 5XX responses than the baseline: \NPetclinicFullFaults per run
on pet-clinic against \NPetclinicBaselineFaults, and
\NPetclinicFullFaultsUnion against \NPetclinicBaselineFaultsUnion over all
runs. This comparison is not between equals. The baseline has no error
mode, while \tool spends the rest of its budget on \mode{BOUNDARY} and
\mode{INVALID} requests, and most of the extra 5XX responses come from them
(Section~\ref{sec:discussion}).

\begin{observation}
\label{obs:faults}
With the same budget, the full configuration produced
\NPetclinicFullFaultsUnion unique 5XX responses on pet-clinic over \NReps
runs, against \NPetclinicBaselineFaultsUnion for the baseline, and
\NGestaohospitalFullFaultsUnion against \NGestaohospitalBaselineFaultsUnion
on gestao-hospital.
\end{observation}

\subsection{\rqref{4}: Cost}
\label{sec:rq4}
Parsing took at most \NParseMsMax~ms per specification
(Table~\ref{tab:parser}). Generating one request took
\NGenUsMedian~$\mu$s at the median over the \NSpecs specifications and
\NGenUsMax~$\mu$s at most, for market, the specification with the most
\texttt{pattern} keywords. Both are small next to one
HTTP round trip to a local API. In requests, \tool makes the nominal phase
shorter, because fewer requests are rejected: on pet-clinic it fell from
\NPetclinicBaselineNominal requests for the baseline to
\NPetclinicFullNominal for the full configuration, which leaves more of the
budget for error requests.

\begin{observation}
\label{obs:cost}
\tool generated a request in \NGenUsMedian~$\mu$s at the median and cut the
nominal phase on pet-clinic from \NPetclinicBaselineNominal to
\NPetclinicFullNominal requests.
\end{observation}
